\documentclass[11pt]{article}
\usepackage[margin=1in]{geometry}
\usepackage{amsmath,amssymb,amsthm}
\usepackage{graphicx}
\usepackage[hidelinks]{hyperref}

\newtheorem{theorem}{Theorem}
\newtheorem{lemma}{Lemma}
\newtheorem{remark}{Remark}

\newcommand{\K}{\mathcal K}
\newcommand{\B}{\mathcal B}
\newcommand{\U}{\mathcal U}
\newcommand{\wh}{\widehat}

\title{BCAP-Free Injectivity of the Calkin Representation}
\author{FA Banach run \texttt{fa\_banach\_001}, agent \texttt{agent\_lane\_05}}
\date{2026-06-16}

\begin{document}
\maketitle

\section*{Status}

This packet records a candidate partial result, likely valid.  It concerns
Remark 1 after Theorem 2.1 of Boedihardjo--Johnson, arXiv:1301.0054, which asks
whether their Calkin representation theorem remains true without the hypothesis
that \(X\) has the bounded compact approximation property (BCAP).

\section*{Source Context}

Boedihardjo and Johnson define, for an infinite-dimensional Banach space \(X\),
an ultrapower space \(\widehat X\) and for \(T\in B(X,Y)\) an induced operator
\(\widehat T\).  Compact operators satisfy \(\widehat K=0\).  Their Theorem 2.1
states that if \(X\) has the \(\lambda\)-BCAP, then
\[
  f:B(X)/K(X)\to B(\widehat X),\qquad \dot T\mapsto \widehat T,
\]
is a norm-one \((\lambda+1)\)-embedding.

\begin{center}
\includegraphics[width=0.98\linewidth]{figures/theorem_context_crop.png}
\end{center}

Immediately after the proof, the authors ask whether the theorem remains true
without the BCAP hypothesis.

\begin{center}
\includegraphics[width=0.98\linewidth]{figures/open_problem_crop.png}
\end{center}

\section*{Result}

\begin{theorem}
Let \(X\) be an infinite-dimensional Banach space.  In the notation of
Boedihardjo--Johnson's Calkin representation, if \(T\in B(X)\) is noncompact,
then \(\widehat T\ne 0\).  Consequently
\[
  f:B(X)/K(X)\to B(\widehat X),\qquad \dot T\mapsto\widehat T,
\]
is injective for every \(X\), without assuming BCAP.
\end{theorem}

\section*{Idea of the Proof}

The BCAP proof uses compact approximants \(S_\alpha\) to produce weak-star-null
functionals on which \(T^*\) stays large.  For injectivity alone, compact
approximants are unnecessary.  If \(T\) is not compact, then \(T(B_X)\) cannot
be uniformly approximated by finite-dimensional subspaces.  Thus there is a
fixed positive lower bound for the quotient norms \(\|Q_E T\|\), where
\(E\subset X\) ranges over finite-dimensional subspaces.  By duality, for each
\(E\) there is a norm-one functional in \(E^\perp\) on which \(T^*\) is still
large.  These functionals form a weak-star-null element of the ultrapower
\(\widehat X\), because the index ultrafilter contains all tails
\(\{E:E\supset E_0\}\).  Applying \(\widehat T\) to this element gives a
nonzero vector.

\begin{proof}
Let \(\Lambda_0\) be the directed set of finite-dimensional subspaces of \(X\),
ordered by inclusion, and let \(\U\) be an ultrafilter containing every tail
\[
  \{E\in\Lambda_0:E\supset E_0\}.
\]
For \(E\in\Lambda_0\), let \(Q_E:X\to X/E\) be the quotient map.

We first record a compactness criterion.  If
\[
  \inf_{E\in\Lambda_0}\|Q_E T\|=0,
\]
then \(T\) is compact.  Indeed, for every \(\varepsilon>0\) there is a
finite-dimensional \(E\subset X\) such that every point of \(T(B_X)\) is within
\(\varepsilon\) of \(E\).  Since \(T(B_X)\) is bounded, it is within
\(\varepsilon\) of a bounded subset of the finite-dimensional space \(E\), and
that bounded subset has a finite \(\varepsilon\)-net.  Hence \(T(B_X)\) is
totally bounded.  This is compactness of \(T\).

Now assume \(T\) is not compact.  Then
\[
  \eta:=\inf_{E\in\Lambda_0}\|Q_E T\|>0.
\]
For every \(E\in\Lambda_0\), duality gives
\[
  \|Q_E T\|=\|(Q_E T)^*\|
  =\|T^*|_{E^\perp}\|.
\]
Choose \(x_E^*\in E^\perp\) with \(\|x_E^*\|=1\) and
\[
  \|T^*x_E^*\|\ge \eta/2.
\]

The bounded family \((x_E^*)_{E\in\Lambda_0}\) defines an ultrapower element
\[
  \xi=(x_E^*)_{E,\U}\in (X^*)^\U.
\]
It lies in \(\widehat X\).  To see this, fix \(x\in X\).  On the tail
\(\{E:E\supset \operatorname{span}\{x\}\}\), which belongs to \(\U\), we have
\(x_E^*(x)=0\).  Hence the weak-star \(\U\)-limit of \(x_E^*\) is zero.
Moreover \(\|\xi\|=1\).

Applying \(\widehat T\) gives
\[
  \widehat T\xi=(T^*x_E^*)_{E,\U}.
\]
Therefore
\[
  \|\widehat T\xi\|
  =\lim_{E,\U}\|T^*x_E^*\|
  \ge \eta/2>0.
\]
Thus \(\widehat T\ne 0\).  Since Boedihardjo--Johnson already observe that
\(\widehat K=0\) for compact \(K\), the induced map on \(B(X)/K(X)\) is
well-defined and injective.
\end{proof}

\section*{What Remains Open}

This proves only faithfulness.  It does not give a constant \(c>0\) such that
\[
  \|\widehat T\|\ge c\|\dot T\|
\]
for all \(T\in B(X)\).  Therefore it does not prove that \(f\) is an isomorphic
embedding without BCAP in the sense of Theorem 2.1.  It also does not settle
the isometry question in Remark 2 for \(L_p(0,1)\), \(p\ne2\).

\section*{Novelty and Literature Check}

The run indexes were searched for arXiv:1301.0054, mean ergodic convergence,
Calkin, BCAP, and Calkin representation phrases.  Existing hits were adjacent
Calkin-algebra problem packets, not this representation question.  A bounded
web search for exact phrases from Remark 1 and Theorem 2.1 did not find a later
answer.  Novelty confidence is moderate: this is a natural observation and may
exist implicitly in operator-theory folklore, but it is not recorded in the run
memory or surfaced by the search.

\section*{Human Review Recommendation}

Reviewers should verify the compactness criterion and the dual identification
\(\|Q_E T\|=\|T^*|_{E^\perp}\|\).  The rest of the argument is then a direct
ultrapower-tail construction.

\begin{thebibliography}{9}

\bibitem{BJ}
M. T. Boedihardjo and W. B. Johnson,
\emph{On Mean Ergodic Convergence in the Calkin Algebras},
arXiv:1301.0054.

\end{thebibliography}

\end{document}
